% Abhakseite „Das kann ich“ – erzeugt von rahmen.py aus den Titeln der Aufgabendateien
\clearpage
\hypertarget{abhaken}{}\einheitenkopf{Das kann ich}
\zweigzeile{Hake ab, was du kannst. Die Nummern sind die Aufgaben auf dem Blatt.}
{\small\setlength{\parskip}{1.5pt}
\abhakgruppe{Kennst du schon}
\abhakzeile{1}{Ich kann Zahlen quadrieren, auch negative Zahlen und Dezimalzahlen.}
\abhakzeile{2}{Ich kann Klammer zuerst und Punkt vor Strich rechnen, auch mit negativen Zahlen.}
\abhakzeile{3}{Ich kann eine lineare Gleichung mit Äquivalenzumformungen lösen.}
\abhakzeile{4}{Ich kann durch Einsetzen prüfen, ob eine Zahl eine Gleichung löst.}
\abhakzeile{5}{Ich kann Quadratwurzeln ziehen und Näherungswerte runden.}
\abhakzeile{6}{Ich kann an der Scheitelpunktform den Scheitelpunkt und die Öffnung einer Parabel ablesen.}
\abhakzeile{7}{Ich kann ausmultiplizieren, ausklammern und binomische Formeln anwenden.}
\abhakzeile{8}{Ich kann Terme ordnen und zusammenfassen.}
\abhakzeile{9}{Ich finde den Fehler beim Quadrieren negativer Zahlen.}
\abhakgruppe{Einheit 1 · Wurzelziehen und Lösbarkeit}
\abhakzeile{10}{Ich erkenne, ob eine Gleichung quadratisch oder linear ist.}
\abhakzeile{11}{Ich erkenne an der rechten Seite von $x^2 = c$, ob es zwei, eine oder keine Lösung gibt.}
\abhakzeile{12}{Ich kann $x^2 = c$ durch Wurzelziehen lösen.}
\abhakzeile{13}{Ich kann $x^2$ erst allein stellen und dann die Wurzel ziehen.}
\abhakzeile{14}{Ich kann eine Gleichung mit quadrierter Klammer lösen.}
\abhakzeile{15}{Ich kann am Graphen ablesen, wie viele Lösungen eine Gleichung hat.}
\abhakzeile{16}{Ich kann Aussagen über die Lösungen quadratischer Gleichungen beurteilen.}
\abhakzeile{17}{Ich finde den Fehler beim Wurzelziehen.}
\abhakzeile{18}{Ich kann begründen, warum $x^2 = c$ zwei, eine oder keine Lösung hat.}
\abhakgruppe{Einheit 2 · Satz vom Nullprodukt}
\abhakzeile{19}{Ich erkenne, ob bei einem Produkt rechts eine Null steht.}
\abhakzeile{20}{Ich kann eine Gleichung in Produktform mit dem Satz vom Nullprodukt lösen.}
\abhakzeile{21}{Ich kann eine Gleichung lösen, indem ich $x$ ausklammere.}
\abhakzeile{22}{Ich finde den Fehler, wenn jemand durch $x$ teilt.}
\abhakzeile{23}{Ich kann begründen, wann der Satz vom Nullprodukt gilt.}
\abhakzeile{24}{Ich kann mit dem Satz vom Nullprodukt die Spannweite eines Brückenbogens bestimmen.}
\abhakgruppe{Einheit 3 · Normalform und p-q-Formel}
\abhakzeile{25}{Ich erkenne, welche Form eine quadratische Gleichung hat und welcher Lösungsweg passt.}
\abhakzeile{26}{Ich kann $p$ und $q$ mit Vorzeichen ablesen.}
\abhakzeile{27}{Ich erkenne, ob die Vorzahl von $x^2$ eins ist.}
\abhakzeile{28}{Ich kann eine Gleichung in Normalform mit der p-q-Formel lösen.}
\abhakzeile{29}{Ich kann eine Gleichung in Normalform mit der p-q-Formel lösen – weiter.}
\abhakzeile{30}{Ich kann eine Gleichung in Normalform mit der p-q-Formel lösen – weiter.}
\abhakzeile{31}{Ich kann unterscheiden, welcher Lösungsweg zu einer Gleichung passt.}
\abhakzeile{32}{Ich kann die Schnittpunkte einer Parabel und einer Geraden berechnen.}
\abhakzeile{33}{Ich finde den Fehler in einer Rechnung mit der p-q-Formel.}
\abhakzeile{34}{Ich kann begründen, wann die p-q-Formel gilt und wie viele Lösungen es gibt.}
\abhakzeile{35}{Ich kann berechnen, wo ein geworfener Ball auf dem Boden aufkommt.}
\abhakgruppe{Einheit 4 · Sachaufgaben}
\abhakzeile{36}{Ich erkenne, welche Lösung zur Frage passt.}
\abhakzeile{37}{Ich kann ein Zahlenrätsel mit einer quadratischen Gleichung lösen.}
\abhakzeile{38}{Ich kann eine Sachaufgabe zum Rechteck mit einer quadratischen Gleichung lösen.}
\abhakzeile{39}{Ich finde den Fehler in einer Sachaufgabe.}
\abhakzeile{40}{Ich kann begründen, wann eine negative Lösung zur Frage passt.}
}
